\subsection{Canonical chart orders and their differential graphs}
\label{sec:geometry-global-chart-orders}

The function-order comparison in
Section~\ref{sec:geometry-witt-additive-lattice} admits a
direct descent proof. In particular a nonflat dominant
change of the original chart does not require a new
estimate of its torsion order.

\begin{theorem}[Completion of a split quasi-finite group]
\label{thm:geometry-canonical-chart-order}
Let \(A\) be a normal noetherian domain with fraction field \(K\),
and \(\mathscr G/\operatorname{Spec}A\) a separated, flat,
quasi-finite group scheme of finite presentation with
specified generic fibre the constant finite commutative group
\(T_K\). Define
\[
 M_A=\{f\in A^T:\text{ for every }t\in T,\ 
          (x\mapsto f(t+x))\text{ extends regularly on }\mathscr G\}.
                                                               \tag{GO1}
\]
The extensions refer to the prescribed generic labels.
Then \(M_A\) is a finite locally free Hopf order in \(K^T\).
The evaluation embedding
\(\mathscr G\to\operatorname{Spec}M_A\) is an open immersion,
and its translates by the specified sections of \(T\) cover
\(\operatorname{Spec}M_A\).
For every injective homomorphism \(A\to A'\) to an integral
domain, with the same generic labels after base change,
\[
 M_{A'}=\operatorname{image}(A'\otimes_A M_A\longrightarrow (A')^T),
                                                               \tag{GO2}
\]
and the indicated map is injective. The definition on the
left uses the actual base-changed group, even when \(A'\)
is a nondiscrete valuation ring or the base change is not flat.
If the original group is the kernel of \([n]\) on a smooth
commutative group scheme, its completed order has the actual
cotangent module \(\omega/n\omega\), where \(\omega\) is the
invariant cotangent bundle of that smooth group. Its evaluation
derivative at the identity is the original differential,
naturally under (GO2).
\end{theorem}

\begin{proof}
First suppose \(A\) is henselian local. The henselian finite-part
theorem gives a finite open and closed subscheme
\(\mathscr H\subset\mathscr G\) containing the closed fibre.
One can obtain this directly by factoring the quasi-finite
separated morphism as an open immersion into a finite scheme
(Zariski's main theorem), and using the decomposition of a
finite algebra over a henselian local ring into its local
factors. A local factor whose closed point belongs to the open
is contained in the open. The finite part is flat; it is a
subgroup since identity, multiplication and inverse preserve
it on the closed fibre. In checking multiplication, the
inverse image of the complementary open and closed part in
the finite scheme \(\mathscr H\times_A\mathscr H\) must be
empty: otherwise its closed fibre would be nonempty.
The same argument checks the other operations.

Put \(H=\mathscr H(K)\subset T\). Quotient by the finite
locally free closed subgroup \(\mathscr H\). The quotient
\(\mathscr Q=\mathscr G/\mathscr H\) is a group algebraic
space and \(\mathscr G\to\mathscr Q\) is an fppf
\(\mathscr H\)-torsor (Stacks, Tag 06PH). In particular
this map is finite locally free, by descent.
The quotient is flat and locally of finite presentation over
\(A\) by faithfully flat descent, and is quasi-compact.
Its identity section is open and closed, as can be checked
after pullback to \(\mathscr G\), where its preimage is
\(\mathscr H\). Its diagonal, obtained by inverse and
difference from that section, is therefore open and closed.
Consequently the quotient is separated and unramified, and
hence \'etale. Its generic fibre is \((T/H)_K\).
The normalization of \(\operatorname{Spec}A\) in this
generic fibre is a disjoint union of \(|T/H|\) copies of
\(\operatorname{Spec}A\): \(A\) is integrally closed.
Zariski's main theorem for quasi-finite separated algebraic
spaces (Stacks, Tags 05W4 and 05W7) thus identifies
\(\mathscr Q\) with
\[
                      \coprod_{\alpha\in T/H} U_\alpha
\]
for open subschemes \(U_\alpha\subset\operatorname{Spec}A\)
containing their respective generic points. The identity open
is the whole base. In applying the normalization version of
the theorem, \(\mathscr Q\) is normal, since it is \'etale
over the normal base; its generic fibre is schematically
dense, since it is flat.

For each coset choose the original generic representative
\(t_\alpha\), choosing \(t_0=0\). Its point in the finite
torsor over \(U_\alpha\) extends uniquely to a section:
on each normal affine open of \(U_\alpha\), its coordinates
are integral over the normal ring and belong to its fraction
field. These sections agree on overlaps. Translation by
this section identifies the corresponding part of
\(\mathscr G\) with \(\mathscr H|_{U_\alpha}\), with the
literal prescribed generic identification \(h\mapsto
t_\alpha+h\).

Now form \(\mathscr B\) as the disjoint union of copies of
\(\mathscr H\) indexed by \(T/H\). Equip it with the
induced group law, namely translation by the section of
\(\mathscr H\) corresponding to
\(t_\alpha+t_\beta-t_{\alpha+\beta}\in H\) when adding in
the two copies. Each such generic section extends over
\(A\), by finiteness and normality as above. The cocycle
identities follow in \(H\) and hence for the sections.
The change of coset representatives gives the canonical
translation isomorphism, so this construction preserves all
generic labels. The group \(\mathscr B\) is finite locally
free and its restriction to the open subschemes just
described is the original \(\mathscr G\). Compatibility of
the group maps can also be checked on the schematically
dense generic fibres of the flat products. All generic
labels extend to sections of \(\mathscr B\), and translating
the identity copy \(\mathscr H\) covers \(\mathscr B\).

Its function algebra in \(A^T\) is precisely
\[
             \{f\in A^T:\ f(t_\alpha+{-})|_H
                              \in\mathcal O(\mathscr H)
                               \text{ for every }\alpha\}.
\]
Every function in this algebra and every translate is
regular on \(\mathscr G\). Conversely (GO1) implies the
displayed conditions by restriction to \(\mathscr H\).
Thus \(\mathcal O(\mathscr B)=M_A\), not just an abstract
isomorphic order. This proves the theorem on a henselian
normal local base.

For a general normal noetherian \(A\), the flatness of
\(\mathscr G/A\) gives injective restriction of its regular
functions to the generic fibre; this can be checked on
affine opens by torsion freeness. The local regularity
condition at every point in (GO1) is equivalent to extension
to a global section, by this uniqueness and the sheaf
property. In particular the set in (GO1) is already an
algebra; taking the algebra it generates, as in
\texttt{btOrder}, does not change it.

Here is the descent check, including the distinction between
flat and nonflat base changes. For a flat extension to an
integral domain \(A_1\), the global sections of
\(\mathscr G\) commute with extension of scalars, by a finite
affine-cover equalizer (intersections are affine since the
morphism is separated over an affine base). The kernel
description
\[
 M_A=\ker\left(A^T\longrightarrow
       \bigoplus_{t\in T} K^T/\Gamma(\mathscr G,\mathcal O),
                   \ f\longmapsto(f(t+{-}))_{t\in T}\right)
\]
therefore commutes with flat base change. The localization
\(K\otimes_A A_1\) injects into the fraction field of \(A_1\);
the functions being tested are in \((A_1)^T\), and the
base-changed global sections already lie in that
localization. Thus the resulting kernel is exactly
\(M_{A_1}\) as defined with the actual fraction field.
Apply this to localization at each prime and its
henselization; these are flat extensions and normal local
domains. The module \(M_A\subset A^T\) is finite because
\(A\) is noetherian. Since it becomes locally free of rank
\(|T|\) on these faithfully flat local extensions, it is
finite projective. Its Hopf operations in generic
coordinates descend from those on the henselizations:
membership in \(M_A\otimes M_A\) and \(M_A\) is checked
faithfully flat locally. The identity and associativity
equations hold generically and hence integrally. The
canonical morphism \(\mathscr G\to\operatorname{Spec}M_A\)
is an open immersion with the translate-cover property by
the same faithfully flat local tests. This constructs
\(\mathscr B\) globally.

Finally let \(A\to A'\) be injective without assuming
flatness. The base change of the just-constructed
\(\mathscr B\) is still finite locally free, has the
prescribed split generic fibre, and contains the actual
\(\mathscr G_{A'}\) as an open with the same translate cover.
The projective \(A'\)-module \(A'\otimes_A M_A\) embeds in
\((A')^T\): its map is injective generically and its domain
is torsion-free. A function satisfying (GO1) over \(A'\)
extends regularly on every translated open, by translating
its extensions on \(\mathscr G_{A'}\). These extensions agree
on overlaps because both equal the specified generic function
and the opens are flat over \(A'\). They glue to a global
function of the finite flat base change of \(\mathscr B\).
The converse holds by restriction and the translation
action. This proves (GO2); no base-change assertion for
global sections of a general nonproper scheme under a
nonflat map has been used.

For the last assertion, the identity section lies in the
open subgroup, so its cotangent module is unchanged by the
completion. The cotangent sequence for the kernel of
multiplication by \(n\) on the smooth group identifies this
module with \(\omega/n\omega\): the map on invariant
differentials is \(n\,1\).
Equivalently, each tangent point over the dual numbers of
\(A/nA\) lies uniquely in the kernel, and evaluation of a
function at that point has first-order term its differential.
This proves both the asserted derivative identification and
its naturality, including in (GO2).
\end{proof}

\begin{corollary}[The original chart function and graph identities]
\label{cor:geometry-original-chart-identity}
Use an original \texttt{ChartPack} and \texttt{chartD}
at an odd correction prime \(p\) unramified in the coefficient
field, with the source's labelled generic torsion and tangent
coordinates and the original full-level corrected image bundles.
For every dominant test-ring extension, its generated
\texttt{funLat} and \texttt{thLat} are the actual induced
finite-field torsion order and the actual differential graph
of Section~\ref{sec:geometry-witt-additive-lattice}.
In particular on the height-one test fields of any common
normal modification with the original bundles pulled back,
the six identities \(wT\subset A_\tau\), (RC1) and (RC2) hold,
using the local additive-lattice theorem there.
\end{corollary}

\begin{proof}
By \texttt{tors\(_1\)} and \texttt{pts\(_1\)} of
\texttt{ChartPack}, its quasi-finite group is exactly the
kernel of \([p]\) with the original labelled generic points.
Theorem~\ref{thm:geometry-canonical-chart-order} identifies
\texttt{btOrder} with the coordinate algebra of its
canonical completion, and (GO2) proves the required generated
base change for \texttt{testOrd}. At a henselian valuation-ring
test point this completion is the induced-coset order:
the finite open part and the prescribed generic translations
characterize it by (GO1). The same construction works at the
higher level \(p^{R_p}\). For the latter level its
character differentials are the original Weil differentials,
by the cubical pairing and the last assertion of the theorem.
Thus there is no freedom to replace the pairing by another
one when applying the local additive-lattice proof.

The unramified coefficient algebra \(\mathcal O_E/p\)
is a product of finite fields. Its action extends from the
original group to the completed order: a pullback of a
function satisfying (GO1) satisfies (GO1), since the action
is a homomorphism with the specified generic labels.
The equations of the coefficient action and the equation
\([p]=0\) hold on the completion by generic density.
For the selected residue field \(F\), the source's
\(pq,sc\) give the projection and inclusion of that product
factor. The idempotents split the completed group scheme
into the product of its actual factors (use \(e,1-e\) for
an idempotent \(e\)). The coordinate restriction to the
chosen factor is surjective, with its pullback along the
projection as a section. In particular \texttt{funLat}
is exactly the function order of this factor in the
prescribed coordinates \(F^n\), and its generated base
change is that same factor of (GO2). This is the literal
\texttt{testOrdW} dictionary, not merely equality of ranks.

All labelled cotangent eigenspaces at this prime are integral
direct summands, because the coefficient algebra is
unramified at \(p\) and its embeddings split over the
original normal base with the prescribed constants.
Over the test ring their characters modulo \(p\) are the
Teichm\"uller lifts with the prescribed residue labels:
the prime-to-\(p\) roots of unity have unique lifts with
that residue and the unramified embedding gives precisely
those lifts. Let \(M\) denote the selected factor's order
and \(M^+\) its augmentation ideal. For a label \(\tau\) the
original \texttt{IsDerAtOne} condition is exactly
\[
         \bar u = d f|_\tau\quad
                     \text{in }L_\tau^\vee/p L_\tau^\vee.
\]
Indeed \texttt{IsDerAtOne} tests the actual value on all
dual-number tangent vectors modulo \(p\); the last assertion
of Theorem~\ref{thm:geometry-canonical-chart-order} makes
this equivalent to the displayed differential congruence.
The character factor's function pulls back to the full
order by the actual projection, as required in
\texttt{thLat}. Evaluating on the original basis
\(b(e(\tau))\) gives precisely its second coordinate.
It follows that the generators in the definition of
\texttt{thLat} already form the graph
\[
        \{(f,u)\in M^+\oplus L_\tau^\vee:
                         d f|_\tau=\bar u\}.
                                                               \tag{GO3}
\]
Any element of the displayed graph admits a full cotangent
lift: take the indicated component and lift its remaining
components using the direct-summand decomposition.
This also checks the converse of the source's existential
definition, not just its inclusion in (GO3).

The first projection in (GO3) is onto \(M^+\), with kernel
\(p L_\tau^\vee\). The augmentation ideal is projective:
it is the complement of the constants in the projective
order \(M\). Hence this sequence splits locally. The graph
of the congruence pulls back to the graph of the actual
base-changed differential even under a nonflat dominant
base change. Explicitly choose a local lift of the map
\(M^+\to L_\tau^\vee/pL_\tau^\vee\) to \(L_\tau^\vee\);
the graph is its graph plus \(0\oplus pL_\tau^\vee\),
and the same description holds after base change.
The surjectivity of the differential onto the labelled
cotangent quotient gives the full second projection.
Its kernel and equivariance give exactly the three
square-ideal/character pins of the torsion-character flag
argument. In particular no assertion about arbitrary
base change of the kernel of a nonprojective image is used.

The original high-level character-differential image
lattice and its dual are those constructed in
Section~\ref{sec:geometry-exact-corrected-flags}.
The quotient and its adjoint give the maps \(q,i\) modulo
\(p^{R_p-2}\) on the original locally free bundles.
They have the exact ranks and form of (HF1), so reduction
modulo \(w\) is exact, and lifts exist on the original
affine charts by projectivity. Their graph lattices
\(\Gamma,\Delta\) are independent of the lifts: changing
\(i\) on \(W\) by \(p^{R_p-2}T\) changes it on
\(w^{-1}W\) by an element of \(T\); changing \(q\) on
\(T\) by \(p^{R_p-2}O^{\rm HT}\) changes it on
\(w^{-1}T\) by an element of \(O^{\rm HT}\).
This also proves that the graphs glue. The projection
conditions of \texttt{Base} follow from (HF1) modulo \(w\)
and the sandwich \(L\subset W\subset w^{-1}L\) and
its cotangent counterpart. In particular the second-coordinate
projection of \(\Gamma\) is the first-coordinate kernel
of \(\Delta\): this is precisely exactness modulo \(w\).
The graphs are locally free via their first projection
sequences. The function order is full after inverting
\(p\), since the torsion group there is \'etale with
split generic fibre and its finite completion is constant;
hence locally it contains a power of \(p\) times the full
function lattice. These verify the base inputs to the
relative Chern calculation.

At every required height-one valuation ring (including
exceptional ones after a common modification), all these
objects are now the actual canonical objects by (GO2),
(GO3), and the locally free pullback just established.
Apply (AL1)--(AL3), quotient generation and the exact
cotangent-sequence argument of the preceding sections.
They give all six original pointwise lattice identities
with the actual fixed generic coordinates. The common
model lemma for relative Chern characters now applies
to these very charts. No postulated compatibility between
unrelated torsion representations has been inserted.
\end{proof}

The corollary uses the already written local additive-lattice
theorem, and thus inherits its proof obligations. It discharges
the separate chart realization/base-change interface; it is
not a substitute for an independent review of that local theorem.
